\documentclass[a4paper,12pt]{article}
\usepackage[utf8]{inputenc}
\usepackage[T1]{fontenc}
\usepackage[ngerman]{babel}
\usepackage{amsmath}
\usepackage{amssymb}
\usepackage{enumitem}
\usepackage{geometry}
\geometry{a4paper, margin=2.5cm}

\title{Einsendeaufgaben -- Aufgabe 4.3\\[0.5em]
\large 61211 Analysis}
\author{}
\date{}

\begin{document}
\maketitle

\section*{Aufgabe 1}

\begin{enumerate}[label=\alph*)]

    \item
    \begin{align*}
        D_1 f(x,y) &= 3x^2 - 2x + 2xy \\
        D_2 f(x,y) &= -2y + x^2 \\
        D_{11} f(x,y) &= 6x - 2 + 2y \\
        D_{12} f(x,y) &= D_{21} f(x,y) = 2x \\
        D_{22} f(x,y) &= -2
    \end{align*}

    $(0,0)$ erfüllt die notwendige Bedingung für Extrema:
    \[
        H(0,0) = \begin{pmatrix} -2 & 0 \\ 0 & -2 \end{pmatrix}
    \]
    \[
        Q(u) = u^T \begin{pmatrix} -2 & 0 \\ 0 & -2 \end{pmatrix} u = -2 u_1^2 - 2 u_2^2
    \]
    Wegen $Q(u) < 0$ für alle $u \in \mathbb{R}^2 \setminus \{0\}$ ist $(0,0)$ ein lokales Maximum. $(0,0)$ ist kein globales Maximum, da $\lim_{x \to \infty} f(x,0) = \infty$.

    Wir bestimmen weitere Kandidaten für Extrema:
    \begin{align*}
        D_1 f(x,y) &= 3x^2 - 2x + 2xy = x(3x - 2 + 2y) = 0 \\
        D_2 f(x,y) &= -2y + x^2 = 0
    \end{align*}
    Umstellung nach $y$:
    \[
        y = \frac{x^2}{2}
    \]
    Einsetzen von $y$ in 1. Gleichung:
    \[
        3x - 2 + 2 \cdot \frac{x^2}{2} = x^2 + 3x - 2 = 0
    \]
    Lösen mit p-q-Formel:
    \[
        x_{1,2} = -\frac{3}{2} \pm \sqrt{\left(\frac{3}{2}\right)^2 + 2} = -\frac{3}{2} \pm \frac{1}{2} \sqrt{17}
    \]

    Wir untersuchen $\Delta(a) := D_{11} f(a) \, D_{22} f(a) - (D_{12} f(a))^2$ für die zwei Punkte $\left(x_1, \frac{x_1^2}{2}\right)$ und $\left(x_2, \frac{x_2^2}{2}\right)$.

    \textbf{Punkt $\left(x_1, \frac{x_1^2}{2}\right)$}
    \[
        x_1 = -\frac{3}{2} + \frac{1}{2}\sqrt{17} > -\frac{3}{2} + \frac{1}{2}\sqrt{16} = \frac{1}{2}
    \]
    \[
        \Delta\left(x_1, \frac{x_1^2}{2}\right) = \underbrace{(\underbrace{6x_1 - 2 + x_1^2}_{>0})(-2)}_{<0} - 4x_1^2
    \]
    also ist $\left(x_1, \frac{x_1^2}{2}\right)$ ein Sattelpunkt.

    \textbf{Punkt $\left(x_2, \frac{x_2^2}{2}\right)$}

    Wir vereinfachen $\Delta\left(x_2, \frac{x_2^2}{2}\right)$:
    \begin{align*}
        \Delta\left(x_2, \frac{x_2^2}{2}\right) &= (6x - 2 + x^2)(-2) - 4x^2 \\
        &= -2(6x - 2 + x^2 + 2x^2) \\
        &= -2(6x - 2 + 3x^2) \\
        &= -2\left(-\frac{6}{2}(3+\sqrt{17}) + \frac{3}{4}(3+\sqrt{17})^2 - 2\right) \\
        &= -\left((3+\sqrt{17})\left(-6 + \frac{3}{2}(3+\sqrt{17})\right) - 4\right)
    \end{align*}
    \[
        -6 + \frac{3}{2}(3+\sqrt{17}) > -6 + \frac{3}{2}(3+4) = -6 + \frac{21}{2} > 1
    \]
    Demnach folgt:
    \[
        \Delta\left(x_2, \frac{x_2^2}{2}\right) = -\Big(\underbrace{(3+\sqrt{17}) \underbrace{\left(-6 + \tfrac{3}{2}(3+\sqrt{17})\right)}_{>1}}_{>4} - 4\Big) < 0
    \]
    Punkt $\left(x_2, \frac{x_2^2}{2}\right)$ ist also auch ein Sattelpunkt.

\end{enumerate}

\end{document}
